% ======================================================================
\section{Physical assumptions}
\label{sec:assumptions}
% ======================================================================

The assumptions below are made, explicitly or implicitly, on the path from
the field to the molecular content. Each names the section where it enters,
says why it is made and where it fails, and gives the direction of the
resulting bias and its measured size where the project has measured it.
Section~\ref{sec:limitations} lists the defects and open issues.

\subsection{Time dependence and equilibrium}
\begin{enumerate}
\item \textbf{Chemical steady state.} The network is integrated until no
  species changes by more than 1\% over a doubling of the integration time, at
  fixed temperature, field, columns and imposed ionisation
  (Section~\ref{sec:chemistry}). No chemical age, advection or time-dependent
  field is represented. The slowest process is \Htwo{} formation, with a time
  scale $1/(R_{\mathrm f}n_{\mathrm H})\approx10^7\,(n_{\mathrm H}/100\,\mathrm{cm^{-3}})^{-1}$\,yr
  for $R_{\mathrm f}\approx3\times10^{-17}$\,cm$^3$\,s$^{-1}$, close to the model's
  coefficient (\dlog{L-33}): $10^7$\,yr at \texttt{hiipdr}'s density and $10^6$\,yr
  at \texttt{pn150k}'s. This is comparable to or longer than the age of an
  H\,\textsc{ii} region and much longer than that of a planetary nebula, so where molecules
  have to form from atomic gas the steady state overstates the molecular
  content, and it describes neither gas just uncovered by an advancing front
  nor molecules inherited from an earlier phase.
\item \textbf{Statistical equilibrium of the \Htwo{} levels}, with the
  excited electronic states treated as instantaneous (radiative lifetimes
  0.5--9\,ns in the data; collisions compete only above
  $\sim10^{14}$\,cm$^{-3}$) (Sections~\ref{sec:balance}, \ref{sec:pump}). The slowest radiative
  relaxation within $X$, the quadrupole decay of the lowest rotational levels
  ($A=2.9\times10^{-11}$\,s$^{-1}$ for $J=2\to0$, about 1100\,yr, in
  \src{dataPDR/\allowbreak{}h2/\allowbreak{}transprob\_X.dat}), is short compared with the dynamical
  times of the modelled objects, and collisions shorten it further.
\item \textbf{Formation closure in the level model}: every molecule removed by
  the Solomon process is replaced by one formed on a grain, and the level
  populations are normalised to one. The formation flux into the levels
  therefore equals the Solomon flux; destruction by other channels and the
  corresponding formation are not in the level balance
  (Section~\ref{sec:balance}). Where the Solomon rate has been shielded away
  and cosmic rays and ion--molecule reactions dominate the destruction,
  formation pumping, and the internal energy it brings, are underestimated
  (Section~\ref{sec:limitations}). With the network's \Htwo\,+\,CRP channels
  alone the other destruction is $0.99\,\zeta$, $10^{-16}$\,s$^{-1}$ at the
  default \kw{pdrZeta}, so this happens where the Solomon rate is shielded
  below a few $10^{-16}$\,s$^{-1}$ (Section~\ref{sec:balance}).
\item \textbf{The coupling is a fixed-point iteration}: shielding columns, level
  populations of other cells, the chemistry's channels in \rtn{ionBalance} and
  the heating terms held in the temperature search are one or two iterations
  behind (Section~\ref{sec:lags}). The lags vanish at a fixed point; because
  the MC field is redrawn at every iteration, the fixed point is statistical,
  and the noise biases cell means through the non-linear rates
  (Section~\ref{sec:mcnoise}).
\item \textbf{Homogeneous cells.} Every quantity is uniform within a cell, and
  the MC field is the cell average. A front thinner than a cell is not
  resolved: on \texttt{hiipdr}'s $31^3$ grid a cell is 0.65\,pc and 0.104\,mag of
  $\AV$ (\dlog{H-026}), so the H/\Htwo{} transition is located only to a cell,
  and an ionisation-front cell spreads over its whole volume the warm, partly
  ionised gas of a layer about 0.18\,pc thick, which then dominates
  [O\,\textsc{i}] 63\,\textmu m (\dlog{L-37}). The shielding columns of the chemistry run to the near face of
  the cell (Section~\ref{sec:avradial}), which under-shields CO in the cells of
  its transition by up to half a cell's column.
\end{enumerate}

\subsection{Radiation field}
\begin{enumerate}\setcounter{enumi}{5}
\item \textbf{Point sources only.} The field is that of the sources of the
  input and their reprocessing; no interstellar radiation field illuminates the
  outer boundary and the cosmic microwave background is not transported
  (Sections~\ref{sec:field}, \ref{sec:Td}). Where the attenuated field of the
  sources falls below the interstellar one, the model's outer gas is more
  molecular and colder, and its grains colder, than the same cloud immersed in
  the interstellar field would be. On \texttt{hiipdr} the source's field falls
  below $\chi=1$ beyond about 0.6 of the outer radius and to a median
  $\chi=0.08$ in the outermost shell; on \texttt{pn150k} it stays above 1.6
  everywhere (\src{output/\allowbreak{}pdrFUV.out} of the fixtures' runs). The grain
  temperature is floored at $T_{\mathrm{CMB}}$ for the chemistry only
  (\dlog{L-30}); the line background of the coolants does include the CMB
  (\dlog{D-377}).
\item \textbf{The MC field is a continuum field.} It contains dust absorption
  and scattering and the gas continuum opacities (including C\,\textsc{i} and
  \Htwo{} photoionisation), but no \Htwo, CO or H\,\textsc{i} Lyman-series
  \emph{line} opacity (Section~\ref{sec:eta}). The photons that the \Htwo{} lines
  absorb remain in the transported field, where they heat dust, eject
  photoelectrons and photoionise C\,\textsc{i}. Behind the H/\Htwo{} transition
  this overstates dust heating and photoelectric emission slightly, since the
  saturated line cores hold a small part of the band's energy (the whole window
  holds 15\% of a Draine-shaped field's), and C\,\textsc{i} photoionisation, which
  happens only in the window, by more: a few tens of per cent at $\AV\sim3$ by
  the project's estimate (Section~\ref{sec:limitations}). Unmeasured.
\item \textbf{Separability of continuum and line attenuation.} The rate seen by
  a molecule is the MC continuum at the line times a line transmission computed
  from column densities along a characteristic path: the straight line to each
  source for the direct beam, and the two-leg path through a scattering site
  for diffuse light. Dust is never applied twice: $\AV=0$ on every ray
  (Section~\ref{sec:rays}). Exact for the direct beam; for the diffuse light it
  assumes that a photon's history is represented by the one site it is
  assigned to.
\item \textbf{Direct beam geometry.} The direct field arrives along the straight
  line from each point source. With several sources the columns and the
  in-cell chord are averages weighted by $r_s^{-2}$ alone, without the sources'
  luminosities or attenuation, applied to the summed direct field
  (Section~\ref{sec:avradial}). Exact for one source, as in all the fixtures;
  for several, the error is unmeasured and has no definite sign.
\item \textbf{Diffuse field geometry.} Its total is the MC's $\chiF$; its
  distribution over twelve fixed directions and over sites along them is the
  formal solution with an isotropic emissivity $\sum\kappa^{\mathrm{sca}}J$ and a
  grey FUV extinction $\RFUV\kappa_V$ (Section~\ref{sec:sites}). The same window
  factor is applied at every site. Measured: the formal solution recovers
  52--130\% of $\chiF$, falling with depth (\dlog{L-4}); the forward-peaked phase
  function ($g=0.29$--0.61 on the fixtures) puts more of the real diffuse field
  on the source side than the weights do.
\item \textbf{Angular quadrature}: twelve HEALPix $N_{\mathrm{side}}=1$
  directions fixed in the grid frame, for the sites and for line escape
  (Sections~\ref{sec:healpix}, \ref{sec:lineRays}). Biased along the ring axis:
  cells nearest the plane $z=0$ hold 5--11\% more \Htwo{} before and across the
  transition than their symmetry-equivalent cells; 48 directions remove this,
  moving the totals by under 1\% (\dlog{L-53}, \dlog{D-410}).
\item \textbf{Spectral shape.} The generic photoreactions assume the Draine
  spectral shape across 6--13.6\,eV (Section~\ref{sec:generic}), while the MC
  field's LW share falls from 0.30 to 0.008 through \texttt{hiipdr}'s PDR
  (\dlog{D-349}); CO takes the ratio of the MC field's LW share to Draine's,
  with the shape inside the window assumed Draine's (Section~\ref{sec:window});
  \Htwo{} uses the MC specific intensity at each line, constant across the
  2\%-wide frequency bin that contains it (Section~\ref{sec:eta}).
\item \textbf{Symmetry.} For a \kw{symmetricXYZ} octant the lower faces are
  mirrors for every walk to the edge of the model; the walks towards the
  sources do not reflect (Section~\ref{sec:raycells}). Valid for a model
  symmetric in the three planes with its sources on them.
\end{enumerate}

\subsection{\texorpdfstring{\Htwo}{H2} microphysics}
\begin{enumerate}\setcounter{enumi}{13}
\item \textbf{Line profile}: Voigt, with Doppler parameter
  $b^2=2kT_e/m_{\mathrm{H_2}}+v_{\mathrm{turb}}^2$ and damping from the
  upper level's total radiative decay rate (to the bound levels and the
  continuum of $X$; the data contain no predissociated levels)
  (Section~\ref{sec:shield}). Damping each line by its own $A$, as Cloudy does,
  would put \texttt{hiipdr}'s H/\Htwo{} crossing 0.21\,mag and \texttt{pn150k}'s
  0.10\,mag shallower (\dlog{H-026}).
\item \textbf{Self-shielding} of each line by the column of its lower level,
  averaged through the cell's own column with the curve of growth; the cell's
  own column is taken over its mean chord along the source direction (direct)
  or over all directions (diffuse) (Section~\ref{sec:ownavg}).
\item \textbf{Line overlap} by the other \Htwo{} lines, as a six-point mean of
  their transmission over the line's absorption; photons absorbed in another
  line's wing are lost, not counted as pumping it; H\,\textsc{i} Lyman lines do
  not shade (Section~\ref{sec:overlap}).
\item \textbf{Trapping}: complete redistribution; a trapped photon re-excites
  the same upper level (on-the-spot); escape probability per line, without
  overlap, averaged over twelve directions (Section~\ref{sec:escape}).
\item \textbf{Colliders}: H (including reactive collisions), He, ortho- and
  para-\Htwo{} and \Hp; no electrons, so the excitation of \Htwo{} in ionised
  gas is underestimated. Missing same-spin pairs (86--99.9\% of them, depending
  on the collider) take g-bar estimates; missing spin-changing pairs get no
  collisional rate; outside the tabulated temperatures (for H, 100--5000\,K) the
  rates are held at the table's end values (Section~\ref{sec:coll}).
\item \textbf{Formation distribution} of \citet{TakahashiUehara2001},
  silicate and carbonaceous surfaces mixed in the proportion of the formation
  rates (Section~\ref{sec:formdist}). Formation through H$^-$ and \Htwo$^+$,
  cosmic-ray excitation and destruction other than Solomon are not in the
  level model.
\end{enumerate}

\subsection{Grains}
\begin{enumerate}\setcounter{enumi}{19}
\item \textbf{\Htwo{} formation} follows
  \citet{CazauxTielens2002,CazauxTielens2004} with the
  \citet{HollenbachMcKee1979} sticking coefficient and 3D-PDR's surface
  energies per material, on the geometric grain area per H nucleus of the
  model's own grains, split into silicate and carbonaceous by the species'
  labels (Sections~\ref{sec:grainarea}, \ref{sec:formation}). On
  \texttt{hiipdr}'s silicates this gives 0.50--0.91 of Cloudy's coefficient,
  mainly through the surface energies and partly through the colder gas; on
  \texttt{pn150k}'s graphite 0.86--1.00 (\dlog{L-33}).
\item \textbf{One grain temperature per cell}: the number-weighted mean over
  species and sizes, the size weights normalised over the whole size grid,
  enters an efficiency that is steep in $T_d$; a size grid wider than a
  species' sizes pulls the mean towards zero and can switch formation off
  (\dlog{H-035}; Section~\ref{sec:Td}).
\item \textbf{Grain temperature} from MOCASSIN's radiative equilibrium with no
  stochastic heating by default (\dlog{L-8}), floored at $T_{\mathrm{CMB}}$ because
  the CMB is not transported (Section~\ref{sec:Td}).
\item \textbf{Extinction ratio} $\RFUV$ weighted by source~1's emitted
  spectrum; it overstates the band extinction of the reddened local field at
  depth, which matters only for the diffuse-site weights (Section~\ref{sec:dust}).
\item \textbf{Photoelectric heating} from the model's grains alone: no PAHs or
  very small grains (\dlog{L-7}), so the heating is lower than PDR models with
  the \citet{BakesTielens1994} prescription would have at the same field, and
  the gas colder (Section~\ref{sec:thermal}).
\end{enumerate}

\subsection{Ownership and coupling}
\begin{enumerate}\setcounter{enumi}{24}
\item \textbf{MOCASSIN owns the ionisation} of H, He, C, O and Mg and the
  electron density: imposed on the chemistry before and after the solve, with
  the photoionisation and recombination rates of these atoms taken from
  MOCASSIN inside the solve (Sections~\ref{sec:linked}, \ref{sec:impose}).
  In particular the \Cp/C ratio, which controls the carbon available to CO, is
  MOCASSIN's, and its C\,\textsc{i} photoionisation has no \Htwo-line shielding.
\item \textbf{MOCASSIN owns the temperatures}: the chemistry and the level
  model run at $T_e$ (clamped to the fitted range of the rate data, while 3D-PDR's
  heating terms are evaluated at the unclamped $T_e$, \dlog{L-30}) and $T_d$
  (Section~\ref{sec:chemistry}). The thermal balance is solved at fixed
  chemistry, with the extra heating terms held fixed during the temperature
  search (Section~\ref{sec:thermal}).
\item \textbf{The chemistry owns the molecular content}, which reaches MOCASSIN
  only as the fraction of each element locked in molecules, as $n(\mathrm{H_2})$
  in the opacity, and through the coolants and heating terms
  (Section~\ref{sec:feedback}).
\item \textbf{Cosmic rays}: a uniform ionisation rate (\kw{pdrZeta}), with no
  attenuation with column (Section~\ref{sec:cr}); their heating, 9.4\,eV per
  \Htwo{} ionisation, is counted in molecular gas only
  \citep{TielensHollenbach1985}.
\end{enumerate}

\subsection{Relation to the standard PDR benchmark}

The PDR code comparison \citep{Rollig2007} fixed, for its benchmark models,
the field in Draine units with $\chi=1.71\,G_0$, $\AV=6.289\times10^{-22}N_{\mathrm H}$,
an FUV attenuation $e^{-3.02\AV}$ wherever a code did not treat it as a
function of wavelength, an \Htwo{} formation coefficient
$3\times10^{-18}T^{1/2}$\,cm$^3$\,s$^{-1}$, a cosmic-ray ionisation rate of
$5\times10^{-17}$\,s$^{-1}$, and photoelectric heating after
\citet{BakesTielens1994}, the only heating process it standardised; the
\Htwo{} shielding was left to each code (their Sect.~4.2.3 and Table~5). The
present model departs from all of these choices by construction: its field is
the transported MC field, scattering included (Section~\ref{sec:split}); its
$\AV/N_{\mathrm H}$, $\RFUV$ and grain area are its own grains'
(Table~\ref{tab:dustfix}); its formation coefficient follows the grains'
temperature (Section~\ref{sec:Td}); $\zeta$ defaults to $10^{-16}$\,s$^{-1}$; its
photoelectric heating comes from its own grains without PAHs; and its \Htwo{}
shielding is the level model's. A benchmark comparison therefore tests the
coupling only after these inputs have been matched, as was done against the
vendored 3D-PDR in \dlog{D-257} (in $\zeta$, composition, $\AV$ conversion and
dust); unmatched, it measures differences of physics rather than of numerics.
